% Generated draft: requires mathematical and source review.
\documentclass[11pt]{article}
\usepackage[margin=25mm]{geometry}
\usepackage{fontspec}
\setmainfont{Latin Modern Roman}
\usepackage{amsmath,amssymb,mathtools,unicode-math}
\setmathfont{Latin Modern Math}
\usepackage{enumitem,needspace,xurl,hyperref}
\setlength{\parindent}{0pt}
\setlength{\parskip}{5pt}
\setlength{\emergencystretch}{3em}
\newcommand{\N}{\mathbb{N}}
\newcommand{\Z}{\mathbb{Z}}
\newcommand{\Q}{\mathbb{Q}}
\newcommand{\R}{\mathbb{R}}
\newcommand{\C}{\mathbb{C}}
\newcommand{\F}{\mathbb{F}}
\newcommand{\A}{\mathbb{A}}
\newcommand{\PP}{\mathbb{P}}
\newcommand{\E}{\mathbb{E}}
\newcommand{\eps}{\varepsilon}
\newcommand{\dd}{\,\mathrm{d}}
\begin{document}
\section*{UnsolvedMath ID: 30001550}
\textbf{Source problem code:} catalogue entry
\subsection{Definitions and mathematical statement}
The supplied record concerns finite words. For a finite word \(w=a_1\cdots a_n\), its length is \(|w|=n\). Let \(\mathcal W\) denote the intended class of words. The record does not specify the alphabet, its allowed size, whether the empty word is admitted, or any further restrictions on \(\mathcal W\). It also supplies no definitions of the parameters \(\tau(w)\) and \(\pi(w)\), or of Assous--Pouzet words. These omissions prevent a fully specified mathematical statement; the following preserves the requested questions as conditional formulations pending those definitions.

\begin{enumerate}
\item Determine whether
\[
\forall w\in\mathcal W,\qquad
|w|>\frac{7}{3}\tau(w)-4
\quad\Longrightarrow\quad
\tau(w)=\pi(w).
\]
The strict inequality and the constant \(4\) are part of the supplied question.

\item Determine the infimum of the constants \(d\) for which there exists a counterexample satisfying
\[
|w|\geq\frac{7}{3}\tau(w)-d.
\]
The meaning of ``counterexample'' and the domain of \(d\) require clarification. In particular, the record does not say whether a counterexample means a word with unequal parameters, a word with \(\tau(w)<\pi(w)\), or a counterexample to the entire implication in the first question. If the intended meaning is \(\tau(w)<\pi(w)\) and \(d\) ranges over \(\R\), the requested quantity has the conditional formulation
\[
\inf\left\{d\in\R:
\exists w\in\mathcal W,\quad
\tau(w)<\pi(w)\ \text{and}\
|w|\geq\frac{7}{3}\tau(w)-d
\right\}.
\]
Neither this interpretation nor attainment of the infimum is established by the supplied record.

\item Determine whether ``such extremal words'' must be Assous--Pouzet words. The record supplies neither a definition of that family nor a precise extremality condition. In particular, it does not specify whether extremality means attainment of the infimum in the preceding question.

\item For each fixed admissible parameter \(\delta\), describe all words in
\[
\mathcal E_{\delta}
=\left\{w\in\mathcal W:
\tau(w)<\pi(w),\quad
|w|\geq\frac{7}{3}\tau(w)-\delta
\right\}.
\]
The domain and any sign restriction on \(\delta\) are unspecified. The strict comparison of the two parameters and the non-strict length bound are preserved exactly.
\end{enumerate}

The record presents these as bundled Ehrenfeucht--Silberger questions. Its supplied dated assessment labels the bundle partially solved and attributes an optimal general \(7/3\) bound to Holub and Nowotka, while reporting that the requested extremal classifications and parameterized descriptions are not all settled. This assessment does not supply definitions or a precise theorem sufficient to resolve the formulation gaps above.
\subsection{Short English statement}
The questions ask how close words with \(\tau(w)<\pi(w)\) can come to a length bound with coefficient \(7/3\), which words are extremal, and how to classify words within a fixed additive allowance. A complete statement requires the missing definitions of the two word parameters, the admissible words and constants, and the Assous--Pouzet family. The supplied partially solved status is dated curation requiring independent review.
\subsection{Sources}
\begin{itemize}
\item Local supplied UnsolvedMath registry record, permanent ID \(30001550\), title ``Extremal Bounds Relating Word Period Parameters,'' source problem code \texttt{catalogue entry}, in the Oberwolfach Reports Open Problems collection. The project identifies \texttt{data/unsolvedmath/problems.json} and the supplementary \texttt{catalog/problems.json} as its local registries; no independent file inspection is asserted here.
\item Supplied source citation: \emph{Mini-Workshop: Combinatorics on Words (2011)}, Oberwolfach Reports, DOI \texttt{10.4171/owr/2010/37}, pp.\ 2195--2244. Original OWR report: \url{https://doi.org/10.4171/owr/2010/37}. The record separately gives proposed year 2010; the relationship between these year labels is not explained.
\item Supplied literature reference, labeled ``Holub--Nowotka solution'': \url{https://oceanrep.geomar.de/20295/1/TheEhrSilProblem.pdf}. No fuller bibliographic citation is supplied.
\item Literature assessment checked at \texttt{2026-08-21}: the supplied curation labels the bundle partially solved, attributes the optimal general \(7/3\) bound to Holub and Nowotka, and says the requested extremal classifications and parameterized descriptions are not all settled. This is supplied dated curation, not an independently verified literature assessment. The background explicitly requests independent verification of time-sensitive status and cited claims.
\item Attribution: the supplied assessment calls these Ehrenfeucht--Silberger questions. No individual proposer is identified. The registry's \texttt{corrected\_verified} statement label is accompanied only by an explanation that the complete statement was preserved while mathematical notation was placed in delimiters; it is not treated here as evidence of mathematical verification. The supplied references have not been visited in preparing this draft.
\end{itemize}
\end{document}
